\section{Gate counts of the complete protocol}
\label{sec:peps_dyadic_circuit_bounds}

The gate counts concern actual occurrences. Two occurrences remain distinct
when their local operators agree. The following consequence of the assembly
bounds is the lifetime count used in the compression argument of
\cite[Section~8]{OpenAI2026PolynomialPEPS}.

\begin{theorem}[Lifetime participation]
    \label{thm:peps_dyadic_circuit_lifetime}
    \leanok
    \uses{def:peps_dyadic_protocol, def:peps_dyadic_protocol_operator}
    \lean{TNLean.PEPS.PairEffect.OriginalCircuit.participationCount_castLayouts}
    \lean{TNLean.PEPS.Approximation.DyadicProtocol.IsDistributionProtocol.participationCount_originalCircuit_le}
    Suppose a supplied protocol satisfies the assembly bounds of
    Definition~\ref{def:peps_dyadic_protocol}, with constants $C,b,k,r$.
    Let $W$ be its complete ordered circuit. Then every party participates in
    at most $b$ nonprivate occurrences of $W$. This count is unchanged when
    equal input and output register layouts are identified.
\end{theorem}
\begin{proof}\leanok
    \uses{thm:peps_original_chronology_composition}
    Equal layout identifications preserve the occurrences. Thus, for each
    party $p$,
    \begin{align}
        n_p(W)=\sum_{t<T}n_p(W_t). \notag
    \end{align}
    The right-hand side is the lifetime sum in the assembly bounds and is
    at most $b$.
\end{proof}

\begin{theorem}[Total nonprivate gate count]
    \label{thm:peps_dyadic_circuit_gate_count}
    \leanok
    \uses{def:peps_dyadic_protocol, def:peps_dyadic_protocol_operator}
    \lean{TNLean.PEPS.Approximation.DyadicProtocol.IsDistributionProtocol.nonprivateCount_originalCircuit_le}
    Suppose a supplied protocol satisfies the assembly bounds of
    Definition~\ref{def:peps_dyadic_protocol}, with constants $C,b,k,r$.
    If it has $T$ changes, its complete ordered circuit has at most $bT$
    nonprivate gate occurrences.
\end{theorem}
\begin{proof}\leanok
    \uses{thm:peps_original_chronology_composition}
    Writing $N(X)$ for the number of nonprivate occurrences in a circuit $X$,
    we have
    \begin{align}
        N(W) & =\sum_{t<T}N(W_t)\leq\sum_{t<T}b=bT. \notag
    \end{align}
\end{proof}

\begin{lemma}[Circuit properties under equal layouts]
    \label{lem:peps_circuit_layout_participants}
    \leanok
    \lean{TNLean.PEPS.PairEffect.OriginalCircuit.castLayouts_rec}
    \lean{TNLean.PEPS.PairEffect.OriginalCircuit.forall_participants_castLayouts}
    Identifying the initial and final register layouts of a circuit with equal
    layouts preserves every property of the circuit. In particular, let $Q$
    be a property of finite sets of parties, and suppose $Q(S_g)$ holds for
    the participating set $S_g$ at every nonprivate occurrence. The same
    assertion holds after the endpoint layouts have been identified.
\end{lemma}
\begin{proof}\leanok
    Substitute the equal endpoint layouts. The circuit and its participating
    sets are unchanged.
\end{proof}

\begin{lemma}[A common stage for the participating set]
    \label{lem:peps_dyadic_circuit_stage_participants}
    \leanok
    \uses{def:peps_dyadic_protocol_operator}
    \lean{TNLean.PEPS.Approximation.DyadicProtocol.exists_stage_participants_originalCircuit}
    Let $W$ be the complete ordered circuit of a supplied protocol with
    $T$ changes. For every nonprivate occurrence $g$ of $W$, there are
    $t<T$ and a nonprivate occurrence $\gamma$ of the change $W_t$ such
    that $S_g=S_\gamma$. Thus all parties in $S_g$ may be referred to
    one common stage occurrence.
\end{lemma}
\begin{proof}\leanok
    \uses{thm:peps_original_chronology_composition, lem:peps_circuit_layout_participants}
    Consider the property that a finite set is the participating set of
    an occurrence at a stage below $T$. Every occurrence of each used
    stage has this property. It is preserved by chronological composition
    and by identifying the equal endpoint layouts.
\end{proof}

\begin{theorem}[Uniform bounds on the original expansions]
    \label{thm:peps_dyadic_circuit_expansions}
    \leanok
    \uses{def:peps_dyadic_protocol, def:peps_dyadic_protocol_operator}
    \lean{TNLean.PEPS.Approximation.DyadicProtocol.IsDistributionProtocol.isExpansionBounded_originalCircuit}
    \lean{TNLean.PEPS.Approximation.DyadicProtocol.IsDistributionProtocol.isMonomialBounded_originalCircuit}
    Suppose a supplied protocol satisfies the assembly bounds of
    Definition~\ref{def:peps_dyadic_protocol}, with constants $C,b,k,r$.
    In its complete ordered circuit, each nonprivate gate has at most
    $\lfloor C L^k\rfloor$ monomials. Each monomial contains at most
    $r$ pair effects, and the sum of the absolute values of the coefficients
    at each gate is at most $C L^k$.
\end{theorem}
\begin{proof}\leanok
    \uses{thm:peps_original_chronology_composition}
    Each change satisfies the two uniform bounds. Chronological composition
    preserves them, and identifying the equal endpoint register layouts
    changes no expansion.
\end{proof}
